\documentclass[12pt, a4paper]{article}

\usepackage[utf8]{inputenc}
\usepackage[T1]{fontenc}
\usepackage[italian]{babel}
\usepackage{geometry}
\geometry{
    a4paper,
    total={170mm,257mm},
    left=20mm,
    top=20mm,
}
\usepackage{array}
\usepackage{longtable} % Necessario per le tabelle multipagina

\begin{document}

\subsection*{Project Scoping Meeting Agenda}

\noindent
\textbf{Dipendente}: \\
\textbf{Data}: 1 maggio \\
\textbf{Orario}: 10:00 \\
\textbf{Luogo}: Meeting virtuale (Microsoft Teams)

\subsubsection*{Oggetto}
Primo incontro conoscitivo e di scoping con NexusLog S.p.A. per il progetto ``Shipping-on-the-Air''.

\subsubsection*{Partecipanti}
% [Sezione lasciata intenzionalmente vuota]
\vspace{1.5cm}

\subsubsection*{Svolgimento}
\renewcommand{\arraystretch}{1.5}
\begin{longtable}{|>{\raggedright\arraybackslash}p{0.20\textwidth}|c|c|>{\raggedright\arraybackslash}p{0.42\textwidth}|}
\hline
\multicolumn{3}{|c|}{\textbf{Prima del meeting}} & \multicolumn{1}{c|}{\textbf{Durante il meeting}} \\
\hline
\textbf{Argomento in agenda} & \textbf{Orario} & \textbf{Direttore} & \textbf{Note} \\
\hline
\endfirsthead

\hline
\multicolumn{3}{|c|}{\textbf{Prima del meeting (Continua)}} & \multicolumn{1}{c|}{\textbf{Durante il meeting}} \\
\hline
\textbf{Argomento in agenda} & \textbf{Orario} & \textbf{Direttore} & \textbf{Note} \\
\hline
\endhead

\hline
\endfoot

\hline
\endlastfoot

Introduzione & 10:00 - 10:15 & & NexusLog è un operatore di primissimo piano a livello nazionale. Ci comunicano sin da subito il loro interesse verso l'approccio agile (Scrumban) della nostra startup, ritenuto ideale per gestire l'elevata incertezza tecnologica del progetto. \\
\hline
Scopo del meeting & 10:15 - 10:20 & & Definire il perimetro del progetto ``Shipping-on-the-Air'', chiarire le ambiguità iniziali e valutare la fattibilità tecnica della soluzione software per la gestione dei droni. \\
\hline
Descrizione dello stato corrente e del problema & 10:20 - 11:00 & & NexusLog vanta una rete capillare (5 poli intermodali, 1.200 automezzi, oltre 50.000 spedizioni B2C al giorno), ma l'attuale modello basato esclusivamente sul trasporto su gomma è saturo. Il traffico urbano causa inefficienze sistematiche, ritardi imprevedibili ed elevati costi operativi per l'ultimo miglio. \vspace{0.5em} \newline Inoltre, l'azienda fatica a competere nel nascente settore delle consegne ultra-rapide per merci leggere e non ha un'infrastruttura adeguata per aggredire il mercato del \textit{food delivery}, un segmento in cui vorrebbe entrare con soluzioni proprietarie. \\
\hline
Descrizione dello stato finale & 11:00 - 11:45 & & Il committente vuole rivoluzionare l'ultimo miglio implementando una flotta di droni nello spazio aereo cittadino, garantendo consegne entro 30 minuti. \vspace{0.5em} \newline Viene sottolineato un punto cruciale: nonostante la forte spinta innovativa, NexusLog vuole proteggere la propria \textit{brand reputation} consolidata. Per questo motivo, l'azienda manterrà rigorosamente \textit{in-house} la gestione operativa e commerciale del servizio, affidando in \textit{outsourcing} esclusivamente lo sviluppo dell'architettura software e tecnologica. \\
\hline
Discussione dei requisiti & 11:45 - 12:30 & & I macro requisiti architetturali individuati per l'orchestrazione del sistema sono:
\begin{itemize}
    \item \textbf{Gestione Ordini (\textit{Delivery \& Order}):} Ricezione e orchestrazione del ciclo di vita.
    \item \textbf{Logistica e Flotta:} Verifica vincoli pre-volo (batteria, \textit{payload}).
    \item \textbf{Navigazione (\textit{Routing \& GPS}):} Calcolo rotte e aggiramento \textit{no-fly zone}.
    \item \textbf{Tracciamento Live:} Telemetria satellitare in tempo reale.
    \item \textbf{Notifiche \& Pagamenti:} Sistema \textit{event-driven} e gateway integrato.
\end{itemize}
I requisiti sono definiti ad alto livello. Saranno necessari meeting dedicati per ogni sottosistema. \\
\hline
Budget e tempi & 12:30 - 13:00 & & Per mitigare l'esposizione finanziaria, il committente ha stanziato un budget limitato. Il progetto partirà come un \textit{Proof of Concept} (PoC) / \textit{Minimum Viable Product} (MVP). \vspace{0.5em} \newline NexusLog necessita di una prima versione funzionante entro pochi mesi per iniziare i test di volo sul campo nelle aree pilota. \\
\hline
Domande varie & 13:00 - 13:30 & & 
\begin{itemize}
    \item Il \textit{procurement} hardware (droni, stazioni) sarà interamente gestito da NexusLog.
    \item Verranno fornite a breve le API proprietarie per interfacciarsi con la telemetria di bordo.
    \item Sarà necessaria un'integrazione software con i database attuali di NexusLog.
\end{itemize} \\
\hline
Documentazione & 13:30 - 14:00 & & Non essendo chiari i dettagli algoritmici, non è possibile stendere la documentazione definitiva. Le specifiche di dettaglio verranno prodotte nei successivi incontri. \\
\end{longtable}

\subsubsection*{Conditions of Satisfaction (CoS)}
\begin{itemize}
    \item Sviluppo di un \textit{Proof of Concept} (PoC) funzionante per testare l'efficacia del servizio ``Shipping-on-the-Air'' prima di investimenti su larga scala.
    \item Rispetto rigoroso del budget limitato, concentrando le risorse esclusivamente sullo sviluppo software (\textit{outsourcing} tecnologico).
    \item Garanzia di abbattimento dei tempi di consegna per l'ultimo miglio (SLA: consegne B2C e \textit{food delivery} garantite entro 30 minuti).
    \item Preservazione della \textit{brand reputation} di NexusLog mediante un software affidabile, in modo da non impattare negativamente il \textit{core business} logistico tradizionale.
    \item Il sistema dovrà integrare con successo i 6 moduli cardine: gestione ordini, controllo flotta e \textit{payload}, \textit{routing} satellitare, tracciamento live, notifiche e pagamenti.
\end{itemize}

\subsubsection*{Valutazione del Meeting}
\begin{itemize}
    \item[\textbf{[SÌ]}] L'agenda è stata rispettata?
    \item[\textbf{[SÌ]}] Tutti i partecipanti hanno interagito?
    \item[\textbf{[SÌ]}] L'obiettivo del meeting è stato raggiunto? (Sono state definite le macro-aree del PoC).
    \item[\textbf{[SÌ]}] Abbiamo chiarito i prossimi step?
    \item[\textbf{[SÌ]}] Questo meeting è stato utile?
    \item[\textbf{[Note]}] Come potrebbe migliorare il prossimo meeting? \textit{Si eviteranno discussioni generali sul business per concentrarsi esclusivamente sulle specifiche tecniche delle API dei droni e sulla gestione del payload.}
\end{itemize}

\subsubsection*{Pianificazione del prossimo meeting}
\begin{itemize}
    \item \textbf{Obiettivo:} Analisi del primo sottosistema (Gestione Ordini e Controllo Vincoli di Flotta/Payload) e revisione delle API fornite dal produttore dei droni.
    \item \textbf{Data:} 05/05/2026
    \item \textbf{Luogo:} Meeting Virtuale (Microsoft Teams)
    \item \textbf{Partecipanti:} % [Sezione lasciata intenzionalmente vuota]
\end{itemize}

\vspace{1.5cm}
\noindent
\textbf{Firma}: \hrulefill \\ \\
\textbf{Data}: 01/05/2026

\end{document}